% ═══════════════════════════════════════════════════════════════
% LH-08d: Lehrerhinweise Repaso Final + Jahresabschluss-Test
% Spanisch Klasse 7, Sequenz 8
% ═══════════════════════════════════════════════════════════════

\documentclass[11pt, a4paper]{article}

\newif\ifswmodus
\swmodusfalse

\usepackage[utf8]{inputenc}
\usepackage[T1]{fontenc}
\usepackage[ngerman]{babel}
\usepackage{helvet}
\renewcommand{\familydefault}{\sfdefault}

\usepackage[margin=2cm]{geometry}
\usepackage{enumitem}
\usepackage{tabularx}
\usepackage{booktabs}
\usepackage{longtable}

\usepackage{xcolor}
\usepackage{amssymb}
\usepackage{tcolorbox}
\tcbuselibrary{skins}

\usepackage{fancyhdr}
\usepackage{lastpage}
\usepackage{needspace}

% ═══════════════════════════════════════════════════════════════
% FARBEN
% ═══════════════════════════════════════════════════════════════

\definecolor{spanisch-color}{HTML}{c41e3a}
\definecolor{grau-color}{HTML}{888888}
\definecolor{hellgrau-color}{HTML}{f5f5f5}
\definecolor{basis-color}{HTML}{2a7a4b}
\definecolor{standard-color}{HTML}{2c5aa0}
\definecolor{erweiterung-color}{HTML}{7b1fa2}
\definecolor{loesung-color}{HTML}{c62828}

\definecolor{sw-dark}{HTML}{000000}
\definecolor{sw-gray}{HTML}{666666}
\definecolor{sw-white}{HTML}{ffffff}

\ifswmodus
    \colorlet{spanisch}{sw-dark}
    \colorlet{grau}{sw-gray}
    \colorlet{hellgrau}{sw-white}
    \colorlet{basis}{sw-dark}
    \colorlet{standard}{sw-dark}
    \colorlet{erweiterung}{sw-dark}
    \colorlet{loesung}{sw-dark}
\else
    \colorlet{spanisch}{spanisch-color}
    \colorlet{grau}{grau-color}
    \colorlet{hellgrau}{hellgrau-color}
    \colorlet{basis}{basis-color}
    \colorlet{standard}{standard-color}
    \colorlet{erweiterung}{erweiterung-color}
    \colorlet{loesung}{loesung-color}
\fi

\newcommand{\fachfarbe}{spanisch}

% ═══════════════════════════════════════════════════════════════
% VARIABLEN
% ═══════════════════════════════════════════════════════════════

\newcommand{\thema}{Repaso Final + Jahresabschluss-Test}
\newcommand{\sequenz}{08}
\newcommand{\block}{d}
\newcommand{\fach}{Spanisch}
\newcommand{\klasse}{7}
\newcommand{\dauer}{90 Min. (Doppelstunde)}
\newcommand{\lerngruppengroesse}{24}
\newcommand{\datum}{\today}

% ═══════════════════════════════════════════════════════════════
% KOPF- UND FUSSZEILE
% ═══════════════════════════════════════════════════════════════

\pagestyle{fancy}
\fancyhf{}
\renewcommand{\headrulewidth}{0.4pt}
\renewcommand{\footrulewidth}{0.4pt}

\lhead{\fach{} -- Klasse \klasse}
\chead{\textbf{LH-\sequenz\block{} (Lehrerhinweise)}}
\rhead{\datum}

\lfoot{}
\cfoot{}
\rfoot{Seite \thepage\ von \pageref{LastPage}}

% ═══════════════════════════════════════════════════════════════
% BEFEHLE
% ═══════════════════════════════════════════════════════════════

\newcommand{\B}{\textcolor{basis}{\textbf{[B]}}}
\newcommand{\Std}{\textcolor{standard}{\textbf{[S]}}}
\newcommand{\E}{\textcolor{erweiterung}{\textbf{[E]}}}

\newcommand{\loesung}[1]{\textcolor{loesung}{\textbf{#1}}}
\newcommand{\es}[1]{\textcolor{\fachfarbe}{\textbf{#1}}}

\newtcolorbox{kurzplan}{
    colback=hellgrau,
    colframe=\fachfarbe,
    colbacktitle=white,
    coltitle=\fachfarbe,
    boxrule=1pt,
    arc=0pt,
    fonttitle=\bfseries,
    attach boxed title to top left={yshift=-2mm, xshift=5mm},
    boxed title style={boxrule=0pt, colback=white},
    title=Kurzplan
}

\newtcolorbox{differenzierung}{
    colback=white,
    colframe=grau,
    colbacktitle=white,
    coltitle=grau,
    boxrule=0.5pt,
    arc=0pt,
    fonttitle=\bfseries,
    attach boxed title to top left={yshift=-2mm, xshift=5mm},
    boxed title style={boxrule=0pt, colback=white},
    title=Differenzierung
}

\newtcolorbox{fehlerbox}{
    colback=white,
    colframe=loesung,
    colbacktitle=white,
    coltitle=loesung,
    boxrule=0.5pt,
    arc=0pt,
    fonttitle=\bfseries,
    attach boxed title to top left={yshift=-2mm, xshift=5mm},
    boxed title style={boxrule=0pt, colback=white},
    title=Häufige Fehler
}

\newtcolorbox{materialbox}{
    colback=white,
    colframe=\fachfarbe,
    colbacktitle=white,
    coltitle=\fachfarbe,
    boxrule=0.5pt,
    arc=0pt,
    fonttitle=\bfseries,
    attach boxed title to top left={yshift=-2mm, xshift=5mm},
    boxed title style={boxrule=0pt, colback=white},
    title=Material-Checkliste
}

% ═══════════════════════════════════════════════════════════════
% DOKUMENT
% ═══════════════════════════════════════════════════════════════

\begin{document}

% ═══════════════════════════════════════════════════════════════
% KOPFBEREICH
% ═══════════════════════════════════════════════════════════════

\begin{center}
    {\Large\bfseries\textcolor{\fachfarbe}{Lehrerhinweise: \thema}}\\[0.3em]
    {\normalsize LH-\sequenz\block{} | \dauer}
\end{center}

\vspace{10pt}

% ═══════════════════════════════════════════════════════════════
% 1. KURZPLAN
% ═══════════════════════════════════════════════════════════════

\section*{1. Stundenverlauf (Kurzplan)}

\textbf{1. Stunde (45 Min.) -- Wiederholung}

\begin{tabularx}{\textwidth}{|c|l|X|l|}
\hline
\textbf{Zeit} & \textbf{Phase} & \textbf{Inhalt} & \textbf{Material} \\
\hline
0--5 & Einstieg & Jahresrückblick mit Bildkarten (Seq. 1-8) & Bildkarten \\
\hline
5--15 & Wiederholung & Schnellabfrage: \es{ser, tener, -ar}-Verben & Tafel \\
\hline
15--25 & Übung 1 & AB-08d: Reisevokabular + Wetter (Aufg. 1-2) & AB \\
\hline
25--35 & Übung 2 & AB-08d: Pretérito konjugieren (Aufg. 3) & AB \\
\hline
35--40 & Sicherung & Lösungen besprechen, Fehler klären & Tafel \\
\hline
40--45 & Überleitung & Hinweise zur Stundenleistung & -- \\
\hline
\end{tabularx}

\vspace{1em}

\textbf{2. Stunde (45 Min.) -- Stundenleistung}

\begin{tabularx}{\textwidth}{|c|l|X|l|}
\hline
\textbf{Zeit} & \textbf{Phase} & \textbf{Inhalt} & \textbf{Material} \\
\hline
0--5 & Vorbereitung & Ruhe, SL ausgeben, Regeln erklären & SL-08d \\
\hline
5--30 & Test & Stundenleistung (25 Min.) & SL-08d \\
\hline
30--35 & Einsammeln & Tests einsammeln (LRS: +5 Min.) & -- \\
\hline
35--42 & Abschluss & Jahresreflexion, Feedback einholen & -- \\
\hline
42--45 & Ausblick & Vorschau Klasse 8, Verabschiedung & -- \\
\hline
\end{tabularx}

% ═══════════════════════════════════════════════════════════════
% 2. DIFFERENZIERUNG
% ═══════════════════════════════════════════════════════════════

\section*{2. Differenzierung}

\begin{differenzierung}
\textbf{Legende:} \B{} Basis (BR) | \Std{} Standard (MR) | \E{} Erweiterung (GY)

\vspace{0.5em}

\textbf{WICHTIG:} Diese Marker sind NUR für die Lehrkraft. Auf Schülermaterial erscheinen KEINE Niveaumarkierungen!
\end{differenzierung}

\vspace{1em}

\textbf{Arbeitsblatt AB-08d:}

\begin{tabularx}{\textwidth}{|l|X|X|X|}
\hline
& \B{} \textbf{Basis} & \Std{} \textbf{Standard} & \E{} \textbf{Erweiterung} \\
\hline
\textbf{Aufg. 1} & Zuordnung mit Bildunterstützung & Zuordnung selbstständig & + Sätze bilden \\
\hline
\textbf{Aufg. 2} & Wortliste nutzen & Wortliste als Kontrolle & Ohne Wortliste \\
\hline
\textbf{Aufg. 3} & 4 von 6 Verben & Alle 6 Verben & + eigene Sätze \\
\hline
\textbf{Aufg. 4} & 2-3 Sätze genügen & 4 Sätze & 5+ Sätze, komplexer \\
\hline
\end{tabularx}

\vspace{1em}

\textbf{Stundenleistung SL-08d:}

\begin{itemize}
    \item \B{} LRS-SuS: Zeitzugabe +5 Min.
    \item \B{} Aufgabe 4: Satzanfänge auf separatem Blatt
    \item \E{} Zusatzpunkte für sprachliche Eleganz bei Aufg. 4
\end{itemize}

% ═══════════════════════════════════════════════════════════════
% 3. HÄUFIGE FEHLER
% ═══════════════════════════════════════════════════════════════

\section*{3. Häufige Fehler}

\begin{fehlerbox}
\begin{tabularx}{\textwidth}{|l|X|X|}
\hline
\textbf{Fehler} & \textbf{Beispiel} & \textbf{Reaktion} \\
\hline
Pretérito 1.P.Sg. & *\es{viajo} statt \es{viajé} & Akzent betonen: ``é macht Vergangenheit!'' \\
\hline
Pretérito 3.P.Sg. & *\es{viaja} statt \es{viajó} & Akzent auf -ó hinweisen \\
\hline
Wetter-Konstruktion & *\es{Es sol} statt \es{Hace sol} & ``Hace + Nomen'' einprägen \\
\hline
Llueve/Nieva & *\es{Hace lluvia} statt \es{Llueve} & ``Regen und Schnee sind Verben!'' \\
\hline
nosotros im Pretérito & Verwechslung Präsens/Pretérito & Identisch! Kontext nutzen \\
\hline
\end{tabularx}
\end{fehlerbox}

% ═══════════════════════════════════════════════════════════════
% 4. MATERIAL-CHECKLISTE
% ═══════════════════════════════════════════════════════════════

\section*{4. Material-Checkliste}

\begin{materialbox}
\begin{itemize}[label=$\square$]
    \item AB-08d.pdf (\lerngruppengroesse{} Kopien)
    \item ML-08d.pdf (1× für Lehrkraft)
    \item SL-08d.pdf (\lerngruppengroesse{} Kopien für Test)
    \item SL-08d-ML.pdf (1× für Korrektur)
    \item Bildkarten zu Sequenzen 1-8
    \item Tafelmarker / Kreide
    \item Optional: LP-08d.html (Tablets/PCs)
    \item Optional: Konjugationstabelle Pretérito (Folie/Tafel)
\end{itemize}
\end{materialbox}

% ═══════════════════════════════════════════════════════════════
% 5. LÖSUNGEN AB-08d
% ═══════════════════════════════════════════════════════════════

\section*{5. Lösungen AB-08d}

\textbf{Merkkasten: Pretérito indefinido (-ar)}

\begin{tabular}{|l|l|l|}
\hline
\textbf{Person} & \textbf{Endung} & \textbf{viajar} \\
\hline
yo & \loesung{-é} & \loesung{viajé} \\
tú & \loesung{-aste} & \loesung{viajaste} \\
él/ella/usted & \loesung{-ó} & \loesung{viajó} \\
nosotros/-as & \loesung{-amos} & \loesung{viajamos} \\
vosotros/-as & \loesung{-asteis} & \loesung{viajasteis} \\
ellos/ellas/ustedes & \loesung{-aron} & \loesung{viajaron} \\
\hline
\end{tabular}

\vspace{1em}

\textbf{Aufgabe 1: Reisevokabular} (4 Punkte)

\begin{tabular}{ll}
\es{el viaje} & $\rightarrow$ \loesung{die Reise} \\
\es{el avión} & $\rightarrow$ \loesung{das Flugzeug} \\
\es{la playa} & $\rightarrow$ \loesung{der Strand} \\
\es{el hotel} & $\rightarrow$ \loesung{das Hotel} \\
\end{tabular}

\vspace{1em}

\textbf{Aufgabe 2: Wettersätze} (4 Punkte)

\begin{tabular}{ll}
a) Es regnet. & \loesung{Llueve.} \\
b) Es schneit. & \loesung{Nieva.} \\
c) Die Sonne scheint. & \loesung{Hace sol.} \\
d) Es ist kalt. & \loesung{Hace frío.} \\
\end{tabular}

\vspace{1em}

\textbf{Aufgabe 3: Pretérito konjugieren} (6 Punkte)

\begin{tabular}{ll}
a) & \loesung{viajé} \\
b) & \loesung{visitaste} \\
c) & \loesung{nadó} \\
d) & \loesung{bailamos} \\
e) & \loesung{comprasteis} \\
f) & \loesung{tomaron} \\
\end{tabular}

\vspace{1em}

\textbf{Aufgabe 4: Urlaubsbericht} (6 Punkte)

Bewertungskriterien:
\begin{itemize}
    \item Mind. 4 Sätze im Pretérito: 4P
    \item Passender Wortschatz: 1P
    \item Sprachliche Korrektheit: 1P
\end{itemize}

% ═══════════════════════════════════════════════════════════════
% 6. LÖSUNGEN SL-08d (KURZFASSUNG)
% ═══════════════════════════════════════════════════════════════

\section*{6. Lösungen SL-08d (Stundenleistung)}

Siehe ausführliche Musterlösung: \textbf{SL-08d-ML.pdf}

\textbf{Kurzübersicht:}

\begin{tabular}{|l|l|l|}
\hline
\textbf{Aufgabe} & \textbf{Punkte} & \textbf{Lösungen} \\
\hline
1: Reisevokabular & 4P & el tren = Zug, el avión = Flugzeug, etc. \\
\hline
2: Wettersätze & 4P & Hace calor, Llueve, Hace frío, Hace sol \\
\hline
3: Pretérito & 6P & viajé, visitaste, nadó, bailamos, comprasteis, tomaron \\
\hline
4: Urlaubsbericht & 6P & Individuelle Lösung (Kriterien beachten) \\
\hline
\end{tabular}

% ═══════════════════════════════════════════════════════════════
% 7. NOTENSCHLÜSSEL
% ═══════════════════════════════════════════════════════════════

\section*{7. Notenschlüssel (Sek I, 14 NP)}

\begin{tabular}{|c|c|c|c|c|c|c|c|c|c|c|c|c|c|c|}
\hline
\textbf{NP} & 14 & 13 & 12 & 11 & 10 & 9 & 8 & 7 & 6 & 5 & 4 & 3 & 2 & 1 \\
\hline
\textbf{\%} & 100 & 96 & 91 & 86 & 80 & 73 & 67 & 60 & 53 & 47 & 40 & 33 & 27 & 20 \\
\hline
\textbf{ab P.} & 20 & 19 & 18 & 17 & 16 & 15 & 13 & 12 & 11 & 9 & 8 & 7 & 5 & 4 \\
\hline
\end{tabular}

% ═══════════════════════════════════════════════════════════════
% 8. REFLEXIONSNOTIZEN
% ═══════════════════════════════════════════════════════════════

\section*{8. Reflexionsnotizen (nach Durchführung)}

\begin{tabularx}{\textwidth}{|l|X|}
\hline
\textbf{Was lief gut?} & \\[2em]
\hline
\textbf{Was lief nicht?} & \\[2em]
\hline
\textbf{Zeitmanagement} & \\[2em]
\hline
\textbf{Für nächstes Jahr} & \\[2em]
\hline
\end{tabularx}

\end{document}
